% SPDX-License-Identifier: CC-BY-4.0
\documentclass[11pt,letterpaper]{article}

\usepackage[T1]{fontenc}
\usepackage{amsmath,amssymb}
\usepackage{tikz}
\usepackage{pgfplots}
\pgfplotsset{compat=1.18}
\usepackage{algorithm}
\usepackage{algpseudocode}
\usepackage{booktabs}
\usepackage{enumitem}
\usepackage[hidelinks]{hyperref}

\newcommand{\N}{\mathbb{N}}
\newcommand{\Z}{\mathbb{Z}}
\newcommand{\Q}{\mathbb{Q}}
\newcommand{\R}{\mathbb{R}}
\newcommand{\C}{\mathbb{C}}
\newcommand{\Map}[3]{#1 \colon #2 \to #3}
\newcommand{\Mod}[1]{#1\text{-}\mathrm{Mod}}
\newcommand{\GroupHom}{\mathrm{Hom}_{\mathbf{Grp}}}
\newcommand{\RingHom}{\mathrm{Hom}_{\mathbf{Ring}}}

\newtheorem{theorem}{Theorem}
\newtheorem{lemma}[theorem]{Lemma}
\newtheorem{corollary}[theorem]{Corollary}
\newtheorem{definition}[theorem]{Definition}
\newtheorem{remark}[theorem]{Remark}
\newenvironment{proof}{\noindent\textit{Proof.}}{\hfill\(\square\)}

\title{Reflection at the negative unit multiplier}
\author{ascendnoosphere}
\date{\today}

\begin{document}
\maketitle

\begin{abstract}
For the multiplier \(-1\), the affine state at an even time is the seed and at
an odd time is the reflection of the seed through the offset. The two-step
recurrence repeats those values, with times \(0\) and \(1\) as base cases.
Consequently the signed divisor meets the affine state at time \(0\), at time
\(1\), or never. A return quotient exists only in the first two cases. If
neither value is divisible, the raw orbit alternates and never divides, so
that alternation is a trap of the undivided branch and not a cycle of the
return map.
\end{abstract}

\section{Affine state and return}

Throughout, \(q\), \(c\), \(n\), and \(d\) denote integers and \(j\) denotes a
natural number, including zero. The offset sum is defined by \(T_0(q)=0\) and
\[
T_{j+1}(q)=T_j(q)+q^j.
\]
The two evaluations used below are \(T_{2k}(-1)=0\) and \(T_{2k+1}(-1)=1\)
for every \(k\in\N\).

\begin{definition}
\label{def:affine-state}
The affine state at time \(j\) is
\[
A_j(q,c,n)=q^j n+c\,T_j(q).
\]
\end{definition}

A parameter record consists of a multiplier \(q\), an offset \(c\), and a
divisor \(d\) with \(|d|\ge 2\). In particular \(d\neq 0\). The raw step is
\[
S(n)=\begin{cases}
n/d, & d\mid n,\\
qn+c, & d\nmid n.
\end{cases}
\]
Write \(S^{0}(n)=n\) and \(S^{j+1}(n)=S(S^{j}(n))\). Division of an exact
multiple is the integer quotient: if \(d\mid m\) and \(m=dk\), then
\(m/d=k\), hence \(m=d\cdot(m/d)\). This uses \(d\neq 0\).

\begin{definition}
\label{def:first-hit}
Time \(j\) is a first hit of the seed \(n\), written
\(\operatorname{Hit}(p,n,j)\), when \(d\mid A_j(q,c,n)\) and
\(d\nmid A_i(q,c,n)\) for every natural number \(i<j\).
\end{definition}

When \(j=0\), the second requirement quantifies over the empty set of earlier
times.

\begin{definition}
\label{def:return}
An integer \(r\) is a return of \(n\), written
\(\operatorname{Ret}(p,n,r)\), when there exists \(j\in\N\) such that
\(\operatorname{Hit}(p,n,j)\) and \(A_j(q,c,n)=dr\). The pair \((n,r)\) is a
two-cycle of the return map when \(\operatorname{Ret}(p,n,r)\) and
\(\operatorname{Ret}(p,r,n)\).
\end{definition}

The return map is defined only at seeds that have a first hit. An undivided
orbit is not in its domain.

\section{Powers of minus one}

\begin{lemma}
\label{lem:pow-even}
For every \(k\in\N\), \((-1)^{2k}=1\).
\end{lemma}

\begin{proof}
The exponent identity \(a^{2k}=(a^2)^k\) gives
\((-1)^{2k}=((-1)^2)^k\). Proceed by induction on \(k\). If \(k=0\), then
\(((-1)^2)^0=1\). If the claim holds at \(k\), then
\[
((-1)^2)^{k+1}=((-1)^2)^k\cdot(-1)^2=1\cdot(-1)^2.
\]
The square is \((-1)^2=(-1)\cdot(-1)=1\), so the product is \(1\).
\end{proof}

\begin{lemma}
\label{lem:pow-odd}
For every \(k\in\N\), \((-1)^{2k+1}=-1\).
\end{lemma}

\begin{proof}
The successor rule \(a^{m+1}=a^m\cdot a\) at \(m=2k\) gives
\[
(-1)^{2k+1}=(-1)^{2k}\cdot(-1).
\]
Lemma~\ref{lem:pow-even} replaces the even power by \(1\), and
\(1\cdot(-1)=-1\).
\end{proof}

\begin{lemma}
\label{lem:parity}
For every \(j\in\N\), either \(j=2k\) for some \(k\in\N\), or \(j=2k+1\) for
some \(k\in\N\).
\end{lemma}

\begin{proof}
Induce on \(j\). The time \(0\) equals \(2\cdot 0\). If \(j=2k\), then
\(j+1=2k+1\). If \(j=2k+1\), then \(j+1=2k+2=2(k+1)\).
\end{proof}

\section{The closed form}

\begin{theorem}
\label{thm:reflection-affine}
Let \(C,n\in\Z\) and let \(j\in\N\). If \(j=2k\) for some \(k\in\N\), then
\[
A_j(-1,C,n)=n.
\]
If \(j=2k+1\) for some \(k\in\N\), then
\[
A_j(-1,C,n)=C-n.
\]
\end{theorem}

\begin{proof}
Suppose first that \(j=2k\). Definition~\ref{def:affine-state} expands the
state as
\[
A_{2k}(-1,C,n)=(-1)^{2k}n+C\,T_{2k}(-1).
\]
The even offset sum is \(T_{2k}(-1)=0\), and Lemma~\ref{lem:pow-even} gives
\((-1)^{2k}=1\). Therefore the state equals \(1\cdot n+C\cdot 0=n\).

Suppose instead that \(j=2k+1\). The same expansion gives
\[
A_{2k+1}(-1,C,n)=(-1)^{2k+1}n+C\,T_{2k+1}(-1).
\]
The odd offset sum is \(T_{2k+1}(-1)=1\), and Lemma~\ref{lem:pow-odd} gives
\((-1)^{2k+1}=-1\). Therefore the state equals \(-n+C=C-n\).
\end{proof}

The two base times are the cases \(k=0\). They are recorded separately
because the repetition argument below takes them as its starting values.

\begin{lemma}
\label{lem:time-zero}
For all \(C,n\in\Z\), \(A_0(-1,C,n)=n\).
\end{lemma}

\begin{proof}
Here \(0=2\cdot 0\), so \(T_0(-1)=T_{2\cdot 0}(-1)=0\) and
\((-1)^0=(-1)^{2\cdot 0}=1\) by Lemma~\ref{lem:pow-even}. Hence
\[
A_0(-1,C,n)=1\cdot n+C\cdot 0=n.
\]
\end{proof}

\begin{lemma}
\label{lem:time-one}
For all \(C,n\in\Z\), \(A_1(-1,C,n)=C-n\).
\end{lemma}

\begin{proof}
Here \(1=2\cdot 0+1\), so \(T_1(-1)=1\) and \((-1)^1=-1\) by
Lemma~\ref{lem:pow-odd}. Hence
\[
A_1(-1,C,n)=-n+C\cdot 1=C-n.
\]
\end{proof}

\section{Two-time repetition}

The offset sum adds two powers at once:
\[
T_{j+2}(q)=T_j(q)+q^j+q^{j+1}.
\]
For \(q=-1\) those two powers cancel, so the affine state is periodic of
period \(2\).

\begin{lemma}
\label{lem:repeat-two}
For all \(C,n\in\Z\) and every \(j\in\N\),
\[
A_{j+2}(-1,C,n)=A_j(-1,C,n).
\]
\end{lemma}

\begin{proof}
Expand both sides by Definition~\ref{def:affine-state}. The offset identity
replaces \(T_{j+2}(-1)\) by
\[
T_j(-1)+(-1)^j+(-1)^{j+1}.
\]
The successor rule gives \((-1)^{j+1}=(-1)^j\cdot(-1)=-(-1)^j\), so
\[
(-1)^j+(-1)^{j+1}=(-1)^j-(-1)^j=0.
\]
For the leading power, \(j+2=2+j\) and
\((-1)^{2+j}=(-1)^2\cdot(-1)^j\). Since \((-1)^2=1\),
\[
(-1)^{j+2}=(-1)^j.
\]
Therefore
\begin{align*}
A_{j+2}(-1,C,n)
&=(-1)^j n+C\bigl(T_j(-1)+0\bigr)\\
&=(-1)^j n+C\,T_j(-1)\\
&=A_j(-1,C,n).
\end{align*}
The integer \(j\) was arbitrary, so the identity holds at every time,
including the base times \(j=0\) and \(j=1\).
\end{proof}

\begin{lemma}
\label{lem:repeat-bases}
For all \(C,n\in\Z\),
\[
A_{2}(-1,C,n)=A_0(-1,C,n)
\quad\text{and}\quad
A_{3}(-1,C,n)=A_1(-1,C,n).
\]
In particular \(A_2(-1,C,n)=n\) and \(A_3(-1,C,n)=C-n\).
\end{lemma}

\begin{proof}
Lemma~\ref{lem:repeat-two} at \(j=0\) is \(A_{0+2}=A_0\). The same lemma at
\(j=1\) is \(A_{1+2}=A_1\). Lemmas~\ref{lem:time-zero} and~\ref{lem:time-one}
replace the base states by \(n\) and \(C-n\).
\end{proof}

\begin{lemma}
\label{lem:from-bases}
For all \(C,n\in\Z\) and every \(k\in\N\),
\[
A_{2k}(-1,C,n)=A_0(-1,C,n)
\quad\text{and}\quad
A_{2k+1}(-1,C,n)=A_1(-1,C,n).
\]
\end{lemma}

\begin{proof}
Induce on \(k\). If \(k=0\), then \(2k=0\) and \(2k+1=1\), so both claims
are equalities of a state with itself. These are the base cases of the
repetition: time \(0\) and time \(1\).

Suppose both claims hold at \(k\). The index identities
\(2(k+1)=2k+2\) and \(2(k+1)+1=(2k+1)+2\), together with
Lemma~\ref{lem:repeat-two}, give
\[
A_{2(k+1)}(-1,C,n)=A_{2k}(-1,C,n)
\]
and
\[
A_{2(k+1)+1}(-1,C,n)=A_{2k+1}(-1,C,n).
\]
The inductive hypothesis replaces the right-hand sides by \(A_0(-1,C,n)\)
and \(A_1(-1,C,n)\).
\end{proof}

\begin{corollary}
\label{cor:by-repeat}
For all \(C,n\in\Z\) and every \(k\in\N\),
\[
A_{2k}(-1,C,n)=n
\quad\text{and}\quad
A_{2k+1}(-1,C,n)=C-n.
\]
\end{corollary}

\begin{proof}
Lemma~\ref{lem:from-bases} reduces the even time to time \(0\) and the odd
time to time \(1\). Lemmas~\ref{lem:time-zero} and~\ref{lem:time-one} supply
the two base values.
\end{proof}

Corollary~\ref{cor:by-repeat} is the same closed form as
Theorem~\ref{thm:reflection-affine}, reached by repeating the base times
instead of invoking the offset-sum evaluation at a general index. Both
derivations are used below: the theorem evaluates an arbitrary time directly,
and the corollary explains why a later time cannot introduce a third residue.

\section{The three hit clauses}

Fix a parameter record with multiplier \(q=-1\), offset \(C\), and divisor
\(d\) satisfying \(|d|\ge 2\). For every seed \(n\in\Z\) and every time \(j\),
Theorem~\ref{thm:reflection-affine} transports across \(q=-1\):
\[
A_j(-1,C,n)=\begin{cases}
n, & j=2k,\\
C-n, & j=2k+1.
\end{cases}
\]

\begin{theorem}
\label{thm:hit-zero}
\(\operatorname{Hit}(p,n,0)\) if and only if \(d\mid n\). If \(d\mid n\), then
\(n=d\cdot(n/d)\), the raw step is \(S(n)=n/d\), and
\(\operatorname{Hit}(p,n,j)\) is false for every \(j>0\).
\end{theorem}

\begin{proof}
By Lemma~\ref{lem:time-zero}, \(A_0(-1,C,n)=n\). Thus \(d\) divides the time-zero
state if and only if \(d\) divides \(n\). There is no natural number \(i<0\),
so the earlier-time requirement in Definition~\ref{def:first-hit} holds
automatically at time \(0\). This is the equivalence.

Now assume \(d\mid n\). Since \(d\neq 0\) and \(n=dk\) for the witness \(k\) of
divisibility, the quotient rule gives \(n/d=k\), hence
\(n=d\cdot(n/d)\). The raw step on a divisible input is division, so
\(S(n)=n/d\).

Suppose \(j>0\) and, for a contradiction, that time \(j\) is a first hit.
Definition~\ref{def:first-hit} then requires \(d\nmid A_0(-1,C,n)\). But
\(A_0(-1,C,n)=n\) and \(d\mid n\), a contradiction. Thus no positive time is a
first hit. The return at this seed is the time-zero quotient \(n/d\), because
\(A_0(-1,C,n)=n=d\cdot(n/d)\) and time \(0\) is a first hit.
\end{proof}

\begin{theorem}
\label{thm:hit-one}
Assume \(d\nmid n\) and \(d\mid C-n\). Then \(\operatorname{Hit}(p,n,1)\),
\[
A_1(-1,C,n)=C-n=d\cdot\frac{C-n}{d},
\]
the raw steps are \(S(n)=C-n\) and \(S(S(n))=(C-n)/d\), and
\(\operatorname{Hit}(p,n,j)\) is false for every \(j\neq 1\).
\end{theorem}

\begin{proof}
Lemma~\ref{lem:time-one} gives \(A_1(-1,C,n)=C-n\). The hypothesis
\(d\mid C-n\) is therefore divisibility at time \(1\). The only natural number
strictly less than \(1\) is \(0\), and \(A_0(-1,C,n)=n\) is not divisible by
\(d\). Hence time \(1\) is a first hit. The quotient rule, using \(d\neq 0\),
gives \(C-n=d\cdot((C-n)/d)\). Combined with the evaluation of the state,
the return is \((C-n)/d\).

Since \(d\nmid n\), the raw step uses the affine branch:
\[
S(n)=(-1)\cdot n+C=C-n.
\]
The value \(C-n\) is divisible by \(d\), so the next raw step divides:
\[
S(C-n)=(C-n)/d.
\]
Substituting \(S(n)=C-n\) yields \(S(S(n))=(C-n)/d\).

It remains to exclude every time other than \(1\). If time \(0\) were a first
hit, Theorem~\ref{thm:hit-zero} would give \(d\mid n\), contrary to hypothesis.
If \(j\ge 2\) were a first hit, Definition~\ref{def:first-hit} would require
\(d\nmid A_1(-1,C,n)\). But \(A_1(-1,C,n)=C-n\) and \(d\mid C-n\). Thus the
first hit is exactly time \(1\).
\end{proof}

\begin{theorem}
\label{thm:no-hit}
Assume \(d\nmid n\) and \(d\nmid C-n\). Then for every \(j\in\N\), time \(j\)
is not a first hit. For every \(k\in\N\),
\[
A_{2k}(-1,C,n)=n,\qquad A_{2k+1}(-1,C,n)=C-n,
\]
and the raw iterates satisfy
\[
S^{2k}(n)=n,\qquad S^{2k+1}(n)=C-n.
\]
There is no return of \(n\), and \(n\) lies on no two-cycle of the return map.
\end{theorem}

\begin{proof}
The affine evaluations are Corollary~\ref{cor:by-repeat}. To see that they
forbid a first hit, fix \(j\) and apply Lemma~\ref{lem:parity}. If \(j=2k\),
then every later even time repeats the time-zero value, so
\(A_j(-1,C,n)=n\). Divisibility at time \(j\) would mean \(d\mid n\), which
is false. If \(j=2k+1\), then every later odd time repeats the time-one
value, so \(A_j(-1,C,n)=C-n\). Divisibility at time \(j\) would mean
\(d\mid C-n\), which is false. In both cases time \(j\) fails
Definition~\ref{def:first-hit}. Since \(j\) was arbitrary, no return exists:
a return requires some first hit.

Suppose there were a return two-cycle. By Definition~\ref{def:return} it
would supply an integer \(r\) with \(\operatorname{Ret}(p,n,r)\). That already
requires a return of \(n\), which has just been excluded. The alternating
pair is therefore not a cycle of the return map.

The same two residues are the raw orbit. Because \(d\nmid n\), the first raw
step is the affine branch, \(S(n)=C-n\). Because \(d\nmid C-n\), the next
raw step is again the affine branch,
\[
S(C-n)=(-1)\cdot(C-n)+C=n.
\]
Induce on \(k\). For \(k=0\), one has \(S^{0}(n)=n\) and \(S^{1}(n)=S(n)=C-n\).
These are the raw base cases. If \(S^{2k}(n)=n\) and \(S^{2k+1}(n)=C-n\), then
\[
S^{2(k+1)}(n)=S(S^{2k+1}(n))=S(C-n)=n
\]
and
\[
S^{2(k+1)+1}(n)=S(S^{2(k+1)}(n))=S(n)=C-n.
\]
Every raw state is \(n\) or \(C-n\), neither of which is divisible by \(d\).
The orbit remains on the undivided branch forever. That is a trap of the raw
map. It is not a cycle of the return map, whose steps exist only after a
division.
\end{proof}

\begin{theorem}
\label{thm:reflection-hit}
Let \(p\) be a parameter record with multiplier \(-1\), offset \(C\), and
divisor \(d\) satisfying \(|d|\ge 2\), and let \(n\in\Z\). Then:
\begin{enumerate}
\item \(\operatorname{Hit}(p,n,0)\) if and only if \(d\mid n\). In that case
the return is \(n/d\), the raw step equals that quotient, and no later time
is a first hit.
\item If \(d\nmid n\) and \(d\mid C-n\), then the first hit is exactly time
\(1\) and the return is \((C-n)/d\). The raw orbit follows one affine step
to \(C-n\) and then divides.
\item If \(d\) divides neither \(n\) nor \(C-n\), then there is no first hit.
Every even time repeats \(n\) and every odd time repeats \(C-n\), including
the base times \(0\) and \(1\), and the raw orbit alternates between those
two values.
\end{enumerate}
A return exists if and only if \(d\mid n\) or \(d\mid C-n\). Thus a return
exists only in the first two clauses. The third clause is a trap, not a
cycle of the return map.
\end{theorem}

\begin{proof}
The three clauses are Theorems~\ref{thm:hit-zero}, \ref{thm:hit-one},
and~\ref{thm:no-hit}, together with the transport of
Theorem~\ref{thm:reflection-affine} from the multiplier \(-1\) to the record.

Suppose a return exists, and let \(j\) be a first-hit time. By
Lemma~\ref{lem:parity} and Theorem~\ref{thm:reflection-affine}, the state at
time \(j\) is \(n\) or \(C-n\). Divisibility of that state yields
\(d\mid n\) or \(d\mid C-n\).

Conversely, if \(d\mid n\), Theorem~\ref{thm:hit-zero} supplies the return
\(n/d\) at time \(0\). If \(d\mid C-n\) and also \(d\mid n\), the same
time-zero return applies; the first hit is not time \(1\), because time \(0\)
already hits. If \(d\mid C-n\) and \(d\nmid n\), Theorem~\ref{thm:hit-one}
supplies the return \((C-n)/d\) at time \(1\). If neither divisibility holds,
Theorem~\ref{thm:no-hit} says that no return exists. Therefore a return
exists precisely on the complement of the third clause, which is the union
of the situations covered by the first two clauses.
\end{proof}

\begin{remark}
\label{rem:not-after-division}
After the return in the first clause, the raw orbit has divided once. After
the return in the second clause, it has taken one affine step and then
divided. Neither clause describes the raw state at any later time. The third
clause describes every raw time because no division ever occurs.
\end{remark}

\end{document}
